% !TEX root = ../main.tex
%% Chapitre 11 — Analyse asymptotique
%% Programme officiel MPSI/MP2I (BO spécial n°1 du 11/02/2021), deuxième semestre,
%% section « Analyse asymptotique » : a) relations de comparaison (o, O, équivalence)
%% en un point et en l'infini ; b) développements limités, Taylor-Young, opérations,
%% DL usuels, applications ; c) comparaison des suites ; d) exemples de
%% développements asymptotiques, formule de Stirling (admise).
%% Résultats vérifiés par verifications/chapitre_11.py.

\begin{ouverture}[asymptotique]{11}{Analyse asymptotique}{Partie 2 · Deuxième semestre}
  {Équivalents\enspace·\enspace Négligeabilité\enspace·\enspace Développements limités\enspace·\enspace
   Taylor--Young\enspace·\enspace Asymptotes\enspace·\enspace Stirling}
  \objectif{Comparer fonctions et suites : domination $O$, négligeabilité $o$, équivalence $\sim$.}
  \objectif{Calculer un développement limité : Taylor--Young, opérations, composition, primitivation.}
  \objectif{Connaître parfaitement les développements limités usuels en $0$.}
  \objectif{Appliquer les DL : limites, équivalents, tangentes et positions, extremums, asymptotes.}
  \objectif{Mener un développement asymptotique : suites implicites, suites d'intégrales.}
  \objectif{Utiliser la formule de Stirling (admise) et ses conséquences.}
  \prerequis{Limites, continuité, dérivabilité, Rolle et accroissements finis (chapitre 6).}
  \prerequis{Suites réelles : limites, encadrements, théorème de Cesàro (chapitre 5).}
  \prerequis{Primitives, intégrales et croissances comparées (chapitre 4).}
\end{ouverture}

%% =====================================================================
%%  L'ESSENTIEL
%% =====================================================================
\essentiel

\essSection{1}{Comparer : domination, négligeabilité, équivalence}

\begin{definition}[Les trois relations de comparaison]
Soit $a\in\R\cup\ens{-\infty,+\infty}$ et $f,g$ définies au voisinage de $a$,
$g$ ne s'annulant pas au voisinage de $a$ (sauf peut-être en $a$). On dit que :
\begin{itemize}[label={\color{otNavy}\textbullet},leftmargin=1.4em,itemsep=0.6mm,topsep=0.8mm]
  \item $f$ est \textbf{dominée} par $g$, noté $f=O(g)$, si $f/g$ est bornée
    au voisinage de $a$ ;
  \item $f$ est \textbf{négligeable} devant $g$, noté $f=o(g)$, si
    $f/g\xrightarrow[a]{}0$ ;
  \item $f$ est \textbf{équivalente} à $g$, noté $f\underset{a}{\sim}g$, si
    $f/g\xrightarrow[a]{}1$.
\end{itemize}
Ces définitions s'adaptent mot pour mot aux suites, avec $n\to+\infty$.
\end{definition}

\begin{propriete}[Les deux visages de l'équivalence]
$f\underset{a}{\sim}g$ si et seulement si $f=g+o(g)$ au voisinage de $a$ :
être équivalent à $g$, c'est être égal à $g$ \emph{à une erreur négligeable
devant $g$ près}. La relation $\sim$ est une relation d'équivalence ; si
$f\underset{a}{\sim}g$, alors $f$ et $g$ ont le même signe au voisinage de
$a$, et l'une admet une limite en $a$ si et seulement si l'autre aussi, avec
la même limite.
\end{propriete}

\begin{exemple}
En $0$ : $x^{2}=o(x)$, $\sin x\sim x$, $1-\cos x\sim\dfrac{x^{2}}{2}$.
En $+\infty$ : $x=o(x^{2})$, et un polynôme est équivalent à son terme
\emph{dominant} ($3x^{3}+x^{2}\sim3x^{3}$), tandis qu'en $0$ il est
équivalent à son terme \emph{de plus bas degré} ($3x^{3}+x^{2}\sim x^{2}$).
\end{exemple}

\begin{attention}
Ne jamais écrire $f\sim0$ (cela signifierait $f$ nulle au voisinage), et ne
jamais \emph{sommer} des équivalents. Ainsi
$\sqrt{x^{2}+x}\underset{+\infty}{\sim}x$ n'autorise pas à conclure
$\sqrt{x^{2}+x}-x\to0$ : la vraie limite est $\tfrac12$. Pour additionner,
on repasse par les développements limités.
\end{attention}

\essSection{2}{Les règles de calcul}

\begin{propriete}[Opérations sur les équivalents]
Si $f_1\underset{a}{\sim}g_1$ et $f_2\underset{a}{\sim}g_2$, alors
$f_1f_2\underset{a}{\sim}g_1g_2$, $\dfrac{f_1}{f_2}\underset{a}{\sim}\dfrac{g_1}{g_2}$
et $f_1^{\alpha}\underset{a}{\sim}g_1^{\alpha}$ (puissance fixe, fonctions
strictement positives). On peut aussi \emph{changer de variable} : si
$f(y)\underset{b}{\sim}g(y)$ et $\varphi\xrightarrow[a]{}b$, alors
$f\circ\varphi\underset{a}{\sim}g\circ\varphi$. En revanche, la somme et la
composition \emph{à gauche} (par $\exp$, $\ln$, etc.) sont interdites.
\end{propriete}

\begin{exemple}[Composer par l'exponentielle est interdit]
En $+\infty$, $x^{2}+x\sim x^{2}$, mais
$\dfrac{\mathrm{e}^{x^{2}+x}}{\mathrm{e}^{x^{2}}}=\mathrm{e}^{x}\to+\infty$ :
les exponentielles ne sont pas équivalentes. L'équivalence ne se transmet
par $\exp$ que si la \emph{différence} tend vers $0$.
\end{exemple}

\begin{propriete}[Croissances comparées, en langage $o$]
Pour $\alpha,\beta,\gamma>0$ : en $+\infty$,
$(\ln x)^{\beta}=o\bigl(x^{\alpha}\bigr)$ et $x^{\alpha}=o\bigl(\mathrm{e}^{\gamma x}\bigr)$ ;
en $0^{+}$, $\lvert\ln x\rvert^{\beta}=o\bigl(x^{-\alpha}\bigr)$. Pour les suites :
$n^{\alpha}=o(a^{n})$ si $a>1$, $a^{n}=o(n!)$ et $n!=o(n^{n})$.
\end{propriete}

\begin{savoirfaire}[Calculer une limite par équivalents]
Remplacer chaque \emph{facteur} d'un produit ou d'un quotient par un
équivalent plus simple, puis conclure. Exemple :
\[
\frac{\sin 3x}{\tan 5x}\underset{0}{\sim}\frac{3x}{5x}=\frac35
\qquad\text{donc}\qquad
\lim_{x\to0}\frac{\sin 3x}{\tan 5x}=\frac35.
\]
Dès qu'une \emph{somme} apparaît (numérateur $\tan x-x$, etc.), les
équivalents seuls ne suffisent plus : on passe aux développements limités.
\end{savoirfaire}

\essSection{3}{Développements limités : définition et théorème de Taylor--Young}

\begin{definition}[Développement limité d'ordre $n$]
$f$ admet un \textbf{développement limité d'ordre $n$ en $a$} (en abrégé
$\mathrm{DL}_{n}(a)$) s'il existe des réels $c_0,\dots,c_n$ tels que, pour
$h$ au voisinage de $0$,
\[
f(a+h)=c_0+c_1h+c_2h^{2}+\dots+c_nh^{n}+o\bigl(h^{n}\bigr).
\]
Les coefficients, s'ils existent, sont \emph{uniques} ; en tronquant, $f$
admet alors un $\mathrm{DL}_{p}(a)$ pour tout $p\leqslant n$. Le changement
de variable $x=a+h$ ramène toujours l'étude en $0$.
\end{definition}

\begin{propriete}[Ordres 0 et 1 ; parité]
$f$ admet un $\mathrm{DL}_{0}(a)$ si et seulement si $f$ est continue
(ou prolongeable par continuité) en $a$ ; un $\mathrm{DL}_{1}(a)$ si et
seulement si $f$ est dérivable en $a$, et alors $c_0=f(a)$, $c_1=f'(a)$.
Si $f$ est paire (resp. impaire) et admet un DL en $0$, ses coefficients
d'indice impair (resp. pair) sont nuls.
\end{propriete}

\begin{theoreme}[Taylor--Young]
Si $f$ est de classe $\mathcal{C}^{n}$ sur un intervalle $I$ et $a\in I$,
alors $f$ admet un $\mathrm{DL}_{n}(a)$ donné par ses dérivées successives :
\[
f(a+h)=\sum_{k=0}^{n}\frac{f^{(k)}(a)}{k!}\,h^{k}+o\bigl(h^{n}\bigr).
\]
C'est la source de tous les développements usuels ci-dessous.
\end{theoreme}

\begin{attention}
Un $\mathrm{DL}_{2}$ ne garantit pas l'existence de $f''$ : l'implication
« dérivées $\Rightarrow$ DL » ne se renverse qu'à l'ordre $1$. En pratique,
on calcule les DL par \emph{opérations} (section 5), presque jamais en
dérivant $n$ fois.
\end{attention}

\essSection{4}{Les développements limités usuels en $0$}

\begin{propriete}[À connaître parfaitement]
\begin{align*}
\mathrm{e}^{x}&=1+x+\frac{x^{2}}{2!}+\dots+\frac{x^{n}}{n!}+o\bigl(x^{n}\bigr)
&\cos x&=1-\frac{x^{2}}{2!}+\dots+(-1)^{p}\frac{x^{2p}}{(2p)!}+o\bigl(x^{2p+1}\bigr)\\[0.6mm]
\operatorname{ch}x&=1+\frac{x^{2}}{2!}+\dots+\frac{x^{2p}}{(2p)!}+o\bigl(x^{2p+1}\bigr)
&\sin x&=x-\frac{x^{3}}{3!}+\dots+(-1)^{p}\frac{x^{2p+1}}{(2p+1)!}+o\bigl(x^{2p+2}\bigr)\\[0.6mm]
\operatorname{sh}x&=x+\frac{x^{3}}{3!}+\dots+\frac{x^{2p+1}}{(2p+1)!}+o\bigl(x^{2p+2}\bigr)
&\frac{1}{1-x}&=1+x+x^{2}+\dots+x^{n}+o\bigl(x^{n}\bigr)\\[0.6mm]
\ln(1+x)&=x-\frac{x^{2}}{2}+\dots+(-1)^{n-1}\frac{x^{n}}{n}+o\bigl(x^{n}\bigr)
&\arctan x&=x-\frac{x^{3}}{3}+\dots+(-1)^{p}\frac{x^{2p+1}}{2p+1}+o\bigl(x^{2p+2}\bigr)
\end{align*}
\[
(1+x)^{\alpha}=1+\alpha x+\frac{\alpha(\alpha-1)}{2!}x^{2}+\dots
  +\frac{\alpha(\alpha-1)\cdots(\alpha-n+1)}{n!}x^{n}+o\bigl(x^{n}\bigr),
\qquad
\tan x=x+\frac{x^{3}}{3}+o\bigl(x^{4}\bigr).
\]
En remplaçant $x$ par $-x$ : $\dfrac{1}{1+x}=1-x+x^{2}-\dots$ ; les cas
$\alpha=\tfrac12$ et $\alpha=-\tfrac12$ de $(1+x)^{\alpha}$ servent sans cesse :
$\sqrt{1+x}=1+\tfrac{x}{2}-\tfrac{x^{2}}{8}+o\bigl(x^{2}\bigr)$.
\end{propriete}

\noindent
\begin{minipage}[c]{0.55\linewidth}
\begin{exemple}[D'où viennent-ils ?]
$\mathrm{e}^{x}$, $\cos$, $\sin$, $\operatorname{ch}$, $\operatorname{sh}$ et
$(1+x)^{\alpha}$ sortent de Taylor--Young ; $\dfrac{1}{1-x}$ de la somme
géométrique $1+x+\dots+x^{n}+\dfrac{x^{n+1}}{1-x}$ ; $\ln(1+x)$ et
$\arctan x$ s'obtiennent en \emph{primitivant} $\dfrac{1}{1+x}$ et
$\dfrac{1}{1+x^{2}}$ (section 5) ; $\tan x$ en divisant $\sin x$ par
$\cos x$. Savoir les retrouver est aussi important que les connaître.
\end{exemple}
\end{minipage}\hfill
\begin{minipage}[c]{0.41\linewidth}
\centering
\begin{tikzpicture}[scale=0.58,line cap=round]
  \clip (-3.55,-2.45) rectangle (3.55,2.75);
  \draw[otMuted,line width=0.5pt,-{Stealth[length=1.6mm]}] (-3.4,0) -- (3.4,0);
  \draw[otMuted,line width=0.5pt,-{Stealth[length=1.6mm]}] (0,-2.3) -- (0,2.6);
  \draw[otGreen!55!white,line width=0.9pt,dash pattern=on 2.2pt off 1.6pt]
    plot[domain=-2.25:2.25,samples=40] ({\x},{\x});
  \draw[otGreenDark,line width=0.9pt,dash pattern=on 3pt off 1.8pt]
    plot[domain=-3.25:3.25,samples=90] ({\x},{\x-\x*\x*\x/6});
  \draw[otGreen,line width=1pt]
    plot[domain=-3.05:3.05,samples=110] ({\x},{\x-\x*\x*\x/6+pow(\x,5)/120});
  \draw[otNavy,line width=1.3pt]
    plot[domain=-3.35:3.35,samples=120] ({\x},{sin(deg(\x))});
  \node[font=\tiny,text=otGreen!60!black] at (2.75,2.35) {$P_{1}$};
  \node[font=\tiny,text=otGreenDark] at (3.1,-1.3) {$P_{3}$};
  \node[font=\tiny,text=otGreen!40!black] at (2.35,0.9) {$P_{5}$};
  \node[font=\tiny,text=otNavy] at (-2.6,-1.35) {$\sin$};
\end{tikzpicture}\\[0.5mm]
\legende{$\sin$ et ses polynômes de Taylor $P_{1}$, $P_{3}$, $P_{5}$ :
chaque ordre colle de plus en plus près de $0$.}
\end{minipage}

\essSection{5}{Opérations sur les développements limités}

\begin{propriete}[Combinaisons, produits, quotients, composées]
Si $f$ et $g$ admettent des $\mathrm{DL}_{n}(0)$ de parties régulières $P$
et $Q$, alors :
\begin{itemize}[label={\color{otNavy}\textbullet},leftmargin=1.4em,itemsep=0.6mm,topsep=0.8mm]
  \item $\lambda f+\mu g$ admet le $\mathrm{DL}_{n}(0)$ de partie régulière
    $\lambda P+\mu Q$ ;
  \item $fg$ admet le $\mathrm{DL}_{n}(0)$ obtenu en \emph{tronquant} $PQ$ à
    l'ordre $n$ ;
  \item si $g(0)\neq0$, $\dfrac{f}{g}$ admet un $\mathrm{DL}_{n}(0)$ : écrire
    $\dfrac1g=\dfrac{1}{g(0)}\cdot\dfrac{1}{1+u}$ avec $u\to0$, puis
    composer avec $\dfrac{1}{1+u}=1-u+u^{2}-\dots$ ;
  \item si $u(0)=0$, $f\circ u$ se développe en substituant le DL de $u$
    dans celui de $f$ et en tronquant (pratiqué sur des exemples).
\end{itemize}
\end{propriete}

\begin{theoreme}[Primitivation d'un développement limité]
Si $f$, continue sur $I\ni0$, admet le $\mathrm{DL}_{n}(0)$
$f(x)=c_0+c_1x+\dots+c_nx^{n}+o\bigl(x^{n}\bigr)$, alors toute primitive $F$ de $f$
admet le $\mathrm{DL}_{n+1}(0)$
\[
F(x)=F(0)+c_0x+\frac{c_1}{2}x^{2}+\dots+\frac{c_n}{n+1}x^{n+1}+o\bigl(x^{n+1}\bigr) :
\]
on primitive terme à terme. C'est ainsi que $\ln(1+x)$ et $\arctan x$
héritent de leurs DL.
\end{theoreme}

\begin{savoirfaire}[Composer proprement, tronquer systématiquement]
Pour un DL d'ordre $n$ : tout ce qui dépasse $x^{n}$ part dans le $o\bigl(x^{n}\bigr)$,
à chaque étape. Exemple, $\mathrm{e}^{x}\sin x$ à l'ordre $3$ :
\[
\Bigl(1+x+\frac{x^{2}}{2}+\frac{x^{3}}{6}+o\bigl(x^{3}\bigr)\Bigr)
\Bigl(x-\frac{x^{3}}{6}+o\bigl(x^{3}\bigr)\Bigr)
=x+x^{2}+\Bigl(\frac12-\frac16\Bigr)x^{3}+o\bigl(x^{3}\bigr)
=x+x^{2}+\frac{x^{3}}{3}+o\bigl(x^{3}\bigr).
\]
Et pour $\tan x=\dfrac{\sin x}{\cos x}$ :
$\dfrac{1}{\cos x}=\dfrac{1}{1-\frac{x^{2}}{2}+o\bigl(x^{3}\bigr)}
=1+\dfrac{x^{2}}{2}+o\bigl(x^{3}\bigr)$, d'où
$\tan x=x+\dfrac{x^{3}}{3}+o\bigl(x^{4}\bigr)$ (parité).
\end{savoirfaire}

\essSection{6}{Applications et développements asymptotiques}

\noindent
\begin{minipage}[c]{0.55\linewidth}
\begin{savoirfaire}[Tangente, position, extremum]
Si $f(a+h)=c_0+c_1h+c_ph^{p}+o\bigl(h^{p}\bigr)$ avec $c_p\neq0$ ($p\geqslant2$),
la tangente en $a$ est $y=c_0+c_1(x-a)$ et le signe de $c_p h^{p}$ donne la
\emph{position locale} : au-dessus si $c_p>0$ et $p$ pair, au-dessous si
$c_p<0$ et $p$ pair, traversée (inflexion) si $p$ est impair. En un point
critique ($c_1=0$) : $c_2>0$ donne un minimum local, $c_2<0$ un maximum ;
si $c_2=0$, on pousse le DL jusqu'au premier terme non nul.
\end{savoirfaire}
\end{minipage}\hfill
\begin{minipage}[c]{0.41\linewidth}
\centering
\begin{tikzpicture}[scale=0.78,line cap=round]
  \draw[otMuted,line width=0.5pt,-{Stealth[length=1.6mm]}] (-0.4,0) -- (3.6,0);
  \draw[otMuted,line width=0.5pt,-{Stealth[length=1.6mm]}] (0,-0.4) -- (0,4.4);
  \draw[otGreenDark,line width=0.9pt,dash pattern=on 3pt off 1.8pt]
    plot[domain=-0.15:3.45,samples=30] ({\x},{\x+0.5});
  \draw[otNavy,line width=1.3pt]
    plot[domain=-0.3:3.45,samples=90] ({\x},{sqrt(\x*\x+\x+1)});
  \node[font=\tiny,text=otGreenDark,rotate=44] at (1.65,1.75) {$y=x+\tfrac12$};
  \node[font=\tiny,text=otNavy,rotate=40] at (1.2,2.45) {$y=\sqrt{x^{2}+x+1}$};
\end{tikzpicture}\\[0.5mm]
\legende{Asymptote lue sur le développement :
$\sqrt{x^{2}+x+1}=x+\tfrac12+\tfrac{3}{8x}+o\bigl(\tfrac1x\bigr)$.}
\end{minipage}

\begin{savoirfaire}[Asymptote oblique par développement en $1/x$]
En $\pm\infty$, factoriser le terme dominant puis poser $u=\dfrac1x\to0$. Si
$f(x)=ax+b+\dfrac{c}{x}+o\bigl(\dfrac1x\bigr)$, la droite $y=ax+b$ est asymptote, et
le signe de $\dfrac{c}{x}$ place la courbe par rapport à elle de chaque côté.
\end{savoirfaire}

\begin{definition}[Développement asymptotique]
Un \textbf{développement asymptotique} d'une fonction ou d'une suite est une
somme de termes de référence (puissances, logarithmes, exponentielles)
classés du plus grand au plus petit, chacun négligeable devant le précédent,
avec un reste négligeable devant le dernier terme écrit. Exemple :
$H_n=\ln n+\gamma+o(1)$ (chapitre 5). Les DL en sont le cas particulier
« échelle des $h^{k}$ ».
\end{definition}

\begin{theoreme}[Formule de Stirling (admise)]
\[
n!\;\underset{+\infty}{\sim}\;\sqrt{2\pi n}\,\Bigl(\frac{n}{\mathrm{e}}\Bigr)^{\!n},
\qquad\text{soit}\qquad
\ln(n!)=n\ln n-n+\frac12\ln(2\pi n)+o(1).
\]
\end{theoreme}

\begin{exemple}[Le coefficient binomial central]
$\displaystyle\binom{2n}{n}=\frac{(2n)!}{(n!)^{2}}
\underset{+\infty}{\sim}\frac{\sqrt{4\pi n}\,(2n/\mathrm{e})^{2n}}
{2\pi n\,(n/\mathrm{e})^{2n}}=\frac{4^{n}}{\sqrt{\pi n}}$ :
les $4^{n}$ chemins se concentrent, l'exercice 23 le retrouvera même
\emph{sans} Stirling, par les intégrales de Wallis.
\end{exemple}

%% =====================================================================
%%  EXERCICES
%% =====================================================================
\exercices
\rubrique{Applications}

\begin{exercice}[1]{Premiers équivalents}
Justifier soigneusement les équivalents suivants :
\begin{questions}
  \item en $0$ : $\sin x\sim x$, $\tan x\sim x$ et $1-\cos x\sim\dfrac{x^{2}}{2}$ ;
  \item en $0$ : $\ln(1+x)\sim x$ et $\mathrm{e}^{x}-1\sim x$ ;
  \item $3x^{3}+x^{2}\sim3x^{3}$ en $\pm\infty$, mais $3x^{3}+x^{2}\sim x^{2}$ en $0$ ;
  \item $\sqrt{x^{2}+x}\sim x$ en $+\infty$. Qu'en est-il en $-\infty$ ?
\end{questions}
\end{exercice}

\begin{exercice}[1]{Limites par équivalents}
Déterminer les limites suivantes :
\begin{tasks}(2)
  \task $\displaystyle\lim_{x\to0}\frac{\sin 3x}{\tan 5x}$
  \task $\displaystyle\lim_{x\to0}\frac{1-\cos x}{x^{2}}$
  \task $\displaystyle\lim_{x\to+\infty}x\ln\Bigl(1+\frac1x\Bigr)$
  \task $\displaystyle\lim_{x\to0}\frac{\mathrm{e}^{2x}-1}{\sin x}$
\end{tasks}
\end{exercice}

\begin{exercice}[1]{DL usuels : échauffement}
Écrire, en $0$ :
\begin{tasks}(2)
  \task le $\mathrm{DL}_{3}$ de $\mathrm{e}^{x}$
  \task le $\mathrm{DL}_{4}$ de $\cos x$
  \task le $\mathrm{DL}_{3}$ de $\ln(1+x)$
  \task le $\mathrm{DL}_{2}$ de $\sqrt{1+x}$
  \task le $\mathrm{DL}_{3}$ de $\dfrac{1}{1-x}$
  \task le $\mathrm{DL}_{3}$ de $\dfrac{1}{1+x}$
\end{tasks}
\end{exercice}

\begin{exercice}[1]{Produits de développements limités}
\begin{questions}
  \item Déterminer le $\mathrm{DL}_{3}(0)$ de $\mathrm{e}^{x}\sin x$.
  \item Déterminer le $\mathrm{DL}_{4}(0)$ de $\cos^{2}x$, de deux façons :
    par produit, puis via $\cos^{2}x=\dfrac{1+\cos 2x}{2}$.
\end{questions}
\end{exercice}

\begin{exercice}[1]{Quotients et composées}
Déterminer les développements limités en $0$ :
\begin{questions}
  \item $\mathrm{DL}_{3}$ de $\tan x$ (écrire $\dfrac{\sin x}{\cos x}$) ;
  \item $\mathrm{DL}_{2}$ de $\dfrac{1}{1+\mathrm{e}^{x}}$ ;
  \item $\mathrm{DL}_{3}$ de $\ln(\cos x)$.
\end{questions}
\end{exercice}

\begin{exercice}[1]{Primitiver un développement limité}
\begin{questions}
  \item Retrouver le $\mathrm{DL}_{5}(0)$ de $\arctan x$ à partir de celui
    de $\dfrac{1}{1+x^{2}}$.
  \item Retrouver le $\mathrm{DL}_{4}(0)$ de $\ln(1+x)$ à partir de celui
    de $\dfrac{1}{1+x}$.
  \item Déterminer le $\mathrm{DL}_{4}(0)$ de $\ln\bigl(1+x^{2}\bigr)$. Pouvait-on
    composer directement ?
\end{questions}
\end{exercice}

\begin{exercice}[1]{Courbe et tangente}
Pour chacune des fonctions suivantes, donner la tangente en $0$ et la
position locale de la courbe par rapport à elle :
\begin{questions}
  \item $f(x)=\mathrm{e}^{x}$ ;
  \item $g(x)=\ln(1+x)$ ;
  \item $h(x)=\sin x$.
\end{questions}
Pour $f$, montrer que la position est en fait \emph{globale} :
$\mathrm{e}^{x}\geqslant1+x$ pour tout $x\in\R$.
\end{exercice}

\begin{exercice}[1]{Équivalents de suites}
\begin{questions}
  \item Montrer que $\sqrt{n+1}-\sqrt{n}\sim\dfrac{1}{2\sqrt{n}}$.
  \item Donner un équivalent simple de $\sin\dfrac1n$, puis de
    $n^{2}+3n$ et de $\ln\Bigl(1+\dfrac{1}{n^{2}}\Bigr)$.
  \item Retrouver $\displaystyle\lim_{n\to+\infty}\Bigl(1+\frac1n\Bigr)^{\!n}=\mathrm{e}$
    à l'aide d'un développement limité.
\end{questions}
\end{exercice}

\rubrique{Entraînement}

\begin{exercice}[2]{Limites plus délicates}
Déterminer les limites suivantes :
\begin{tasks}(2)
  \task $\displaystyle\lim_{x\to0}\bigl(\cos x\bigr)^{1/x^{2}}$
  \task $\displaystyle\lim_{x\to0}\frac{(1+x)^{1/x}-\mathrm{e}}{x}$
  \task $\displaystyle\lim_{x\to0}\frac{\tan x-x}{x-\sin x}$
  \task $\displaystyle\lim_{x\to+\infty}\Bigl(x-x^{2}\ln\Bigl(1+\frac1x\Bigr)\Bigr)$
\end{tasks}
\end{exercice}

\begin{exercice}[2]{Développements en un point quelconque}
Déterminer :
\begin{questions}
  \item le $\mathrm{DL}_{2}$ de $\ln x$ en $1$ et le $\mathrm{DL}_{3}$ de
    $\mathrm{e}^{x}$ en $1$ ;
  \item le $\mathrm{DL}_{2}$ de $\sqrt{x}$ en $4$ ;
  \item le $\mathrm{DL}_{2}$ de $\sin x$ en $\dfrac{\pi}{6}$.
\end{questions}
\end{exercice}

\begin{exercice}[2]{Prolongement et dérivabilité par DL}
\begin{questions}
  \item Montrer que $f(x)=\dfrac{\mathrm{e}^{x}-1}{x}$ se prolonge par
    continuité en $0$, que ce prolongement est dérivable en $0$, et calculer
    $f'(0)$.
  \item Montrer que $g(x)=\dfrac{\sin x}{x}$ se prolonge en $0$ en une
    fonction dérivable avec $g'(0)=0$, et que $0$ est un maximum local de
    ce prolongement.
\end{questions}
\end{exercice}

\begin{exercice}[2]{Extremums : l'ordre 2… et au-delà}
Pour chacune des fonctions suivantes, $0$ est un point critique ; déterminer
si c'est un extremum local, et sa nature le cas échéant :
\begin{questions}
  \item $f(x)=x^{2}+x^{3}$ ;
  \item $g(x)=x^{3}+x^{4}$ ;
  \item $h(x)=\cos x+\dfrac{x^{2}}{2}$.
\end{questions}
\emph{Pour c), la condition d'ordre 2 ne décide pas : pousser le DL.}
\end{exercice}

\begin{exercice}[2]{Asymptotes par développement}
Soit $f(x)=\sqrt{x^{2}+x+1}$.
\begin{questions}
  \item Montrer que $f(x)=x+\dfrac12+\dfrac{3}{8x}+o\Bigl(\dfrac1x\Bigr)$
    en $+\infty$. En déduire l'asymptote et la position de la courbe.
  \item Mener l'étude analogue en $-\infty$.
\end{questions}
\end{exercice}

\begin{exercice}[2]{La formule de Stirling en action}
En utilisant la formule de Stirling (admise) :
\begin{questions}
  \item montrer que $\displaystyle\binom{2n}{n}\sim\frac{4^{n}}{\sqrt{\pi n}}$ ;
  \item justifier l'écriture
    $\ln(n!)=n\ln n-n+\dfrac12\ln(2\pi n)+o(1)$ ;
  \item en déduire $\bigl(n!\bigr)^{1/n}\sim\dfrac{n}{\mathrm{e}}$.
\end{questions}
\end{exercice}

\begin{exercice}[2]{Le développement de $\bigl(1+1/n\bigr)^{n}$}
On pose $u_n=\Bigl(1+\dfrac1n\Bigr)^{\!n}$.
\begin{questions}
  \item Montrer que
    $\ln u_n=1-\dfrac{1}{2n}+\dfrac{1}{3n^{2}}+o\Bigl(\dfrac{1}{n^{2}}\Bigr)$.
  \item En déduire
    $u_n=\mathrm{e}\Bigl(1-\dfrac{1}{2n}+\dfrac{11}{24n^{2}}+o\Bigl(\dfrac{1}{n^{2}}\Bigr)\Bigr)$,
    puis $\displaystyle\lim_{n\to+\infty}n\bigl(\mathrm{e}-u_n\bigr)=\frac{\mathrm{e}}{2}$.
\end{questions}
\end{exercice}

\begin{exercice}[2]{Une réciproque au microscope}
Pour $y\in\R$, l'équation $x+\ln x=y$ possède une unique solution
$x_y\in\intOO{0}{+\infty}$ (chapitre 6).
\begin{questions}
  \item Montrer que $x_y\xrightarrow[y\to+\infty]{}+\infty$, puis que
    $x_y\underset{+\infty}{\sim}y$.
  \item Montrer que $x_y=y-\ln y+o(\ln y)$ quand $y\to+\infty$.
\end{questions}
\end{exercice}

\begin{exercice}[2]{Les solutions de $\tan x=x$}
\begin{questions}
  \item Montrer que pour tout $n\in\N^{*}$, l'équation $\tan x=x$ possède
    une unique solution $x_n$ dans
    $\Bigl]n\pi-\dfrac{\pi}{2},\,n\pi+\dfrac{\pi}{2}\Bigr[$, et que
    $x_n\in\Bigl]n\pi,\,n\pi+\dfrac{\pi}{2}\Bigr[$.
  \item Montrer que $x_n\sim n\pi$, puis établir le développement
    $x_n=n\pi+\dfrac{\pi}{2}-\dfrac{1}{n\pi}+o\Bigl(\dfrac1n\Bigr)$.
\end{questions}
\end{exercice}

\begin{exercice}[2]{Une suite d'intégrales}
Pour $n\in\N$, on pose $I_n=\displaystyle\int_{0}^{1}\frac{x^{n}}{1+x}\,\mathrm{d}x$.
\begin{questions}
  \item Montrer que $(I_n)$ est décroissante, positive, et que $I_n\to0$.
  \item Établir $I_n+I_{n+1}=\dfrac{1}{n+1}$, et en déduire
    $I_n\sim\dfrac{1}{2n}$.
  \item Montrer que
    $\displaystyle\sum_{k=1}^{n}\frac{(-1)^{k-1}}{k}=\ln2+(-1)^{n-1}I_n$,
    et conclure sur la limite de ces sommes.
\end{questions}
\end{exercice}

\rubrique{Approfondissement}

\begin{exercice}[3]{Moyennes généralisées}
Soit $0<a<b$. Pour $t\neq0$, on pose
$m(t)=\Bigl(\dfrac{a^{t}+b^{t}}{2}\Bigr)^{\!1/t}$.
\begin{questions}
  \item Reconnaître $m(1)$ et $m(-1)$.
  \item Montrer que $m(t)\xrightarrow[t\to0]{}\sqrt{ab}$, la moyenne
    géométrique. \emph{Développer $a^{t}=\mathrm{e}^{t\ln a}$ à l'ordre 1.}
  \item Montrer que $m(t)\xrightarrow[t\to+\infty]{}b$ et
    $m(t)\xrightarrow[t\to-\infty]{}a$.
\end{questions}
\end{exercice}

\begin{exercice}[3]{Héron : la convergence quadratique quantifiée}
Soit $u_0=2$ et $u_{n+1}=\dfrac12\Bigl(u_n+\dfrac{2}{u_n}\Bigr)$ ; on sait
(chapitres 5 et 6) que $u_n\to\sqrt2$ en décroissant, et on pose
$e_n=u_n-\sqrt2$.
\begin{questions}
  \item Montrer que $e_{n+1}=\dfrac{e_n^{2}}{2u_n}\leqslant\dfrac{e_n^{2}}{2\sqrt2}$.
  \item On pose $\rho=\dfrac{e_0}{2\sqrt2}$. Montrer par récurrence que
    $e_n\leqslant2\sqrt2\,\rho^{\,2^{n}}$.
  \item En déduire que le nombre de décimales exactes \emph{double}
    (au moins) à chaque itération, et estimer $e_4$.
\end{questions}
\end{exercice}

\begin{exercice}[3]{Le grand classique $u_{n+1}=\sin u_n$}
Soit $u_0=1$ et $u_{n+1}=\sin u_n$.
\begin{questions}
  \item Montrer que $u_n\to0$ en décroissant, avec $u_n>0$ pour tout $n$.
  \item Montrer que
    $\dfrac{1}{u_{n+1}^{2}}-\dfrac{1}{u_n^{2}}\xrightarrow[n\to+\infty]{}\dfrac13$.
  \item À l'aide du théorème de Cesàro (chapitre 5), en déduire
    $u_n\sim\sqrt{\dfrac{3}{n}}$.
\end{questions}
\end{exercice}

\begin{exercice}[3]{Étude asymptotique complète}
Soit $f(x)=x\,\mathrm{e}^{1/x}$, définie sur $\R^{*}$.
\begin{questions}
  \item Étudier les limites de $f$ en $0^{+}$, $0^{-}$ et $\pm\infty$ ;
    montrer que $f$ se prolonge par continuité à gauche de $0$, et que ce
    prolongement est dérivable à gauche en $0$.
  \item Étudier les variations de $f$ sur chaque intervalle.
  \item Montrer que $f(x)=x+1+\dfrac{1}{2x}+o\Bigl(\dfrac1x\Bigr)$ en
    $\pm\infty$ ; en déduire l'asymptote commune et la position de la
    courbe de part et d'autre. Esquisser l'allure globale.
\end{questions}
\end{exercice}

\begin{exercice}[3]{Wallis, sans Stirling}
On reprend $I_n=\displaystyle\int_{0}^{\pi/2}\sin^{n}t\,\mathrm{d}t$
(chapitre 4) : $(I_n)$ est décroissante, strictement positive, et
$I_n=\dfrac{n-1}{n}\,I_{n-2}$ pour $n\geqslant2$.
\begin{questions}
  \item Montrer que la suite $\bigl(n\,I_n\,I_{n-1}\bigr)_{n\geqslant1}$ est
    constante, égale à $\dfrac{\pi}{2}$.
  \item Montrer que $\dfrac{I_{n+1}}{I_n}\to1$, puis que
    $I_n\sim\sqrt{\dfrac{\pi}{2n}}$.
  \item Établir $I_{2n}=\dfrac{\pi}{2}\cdot\dfrac{1}{4^{n}}\binom{2n}{n}$,
    et en déduire $\displaystyle\binom{2n}{n}\sim\frac{4^{n}}{\sqrt{\pi n}}$
    — sans la formule de Stirling.
\end{questions}
\end{exercice}

\begin{exercice}[3]{Vers Stirling : encadrer $\ln(n!)$}
\begin{questions}
  \item Montrer que pour tout $k\geqslant1$ :
    $\ln k\leqslant\displaystyle\int_{k}^{k+1}\ln t\,\mathrm{d}t\leqslant\ln(k+1)$.
  \item En sommant, montrer que pour tout $n\geqslant1$ :
    \[
    n\ln n-n+1\;\leqslant\;\ln(n!)\;\leqslant\;n\ln n-n+1+\ln n.
    \]
  \item En déduire $\bigl(n!\bigr)^{1/n}\sim\dfrac{n}{\mathrm{e}}$, et situer ce
    résultat par rapport à la formule de Stirling.
\end{questions}
\end{exercice}

% !TEX root = ../../main.tex
%% Défi du chapitre 11 — La fonction W de Lambert
%% Résultats vérifiés par verifications/defis_11_15.py.

\rubrique{Défi}

\begin{defi}{La fonction W de Lambert}{2 h}
L'équation $x\,\mathrm{e}^{x}=y$ n'a pas de solution « exprimable » à l'aide
des fonctions usuelles ; on en fait donc une fonction nouvelle, la
fonction $W$ de Lambert, omniprésente en physique, en combinatoire et en
informatique. On l'étudie ici entièrement par l'analyse asymptotique.

\partieExo{A}{Définition et dérivée}
\begin{questions}
  \item Étudier $g:x\mapsto x\,\mathrm{e}^{x}$ sur $\intFO{-1}{+\infty}$ et
    montrer qu'elle réalise une bijection de $\intFO{-1}{+\infty}$ sur
    $\intFO{-\frac1{\mathrm{e}}}{+\infty}$. On note $W$ sa réciproque.
    Préciser $W(0)$, $W(\mathrm{e})$ et $W\bigl(-\frac1{\mathrm{e}}\bigr)$.
  \item Montrer que $W$ est continue et strictement croissante, dérivable
    sur $\intOO{-\frac1{\mathrm{e}}}{+\infty}$, avec
    \[
    W'(y)=\frac{1}{\bigl(1+W(y)\bigr)\mathrm{e}^{W(y)}}
    \qquad\text{et}\qquad
    W'(y)=\frac{W(y)}{y\bigl(1+W(y)\bigr)}\ \ (y\neq0).
    \]
    Que vaut $W'(0)$ ? Que se passe-t-il en $-\frac1{\mathrm{e}}$ ?
\end{questions}

\medskip
\begin{center}
\begin{tikzpicture}[scale=1.35,line cap=round]
  \draw[otRule,line width=0.3pt] (-1.2,-1.2) grid (2.3,2.3);
  \draw[otMuted,line width=0.5pt,-{Stealth[length=1.6mm]}] (-1.3,0) -- (2.45,0);
  \draw[otMuted,line width=0.5pt,-{Stealth[length=1.6mm]}] (0,-1.3) -- (0,2.45);
  \draw[otMuted,line width=0.5pt,dash pattern=on 2pt off 1.6pt] (-1.2,-1.2) -- (2.3,2.3);
  \draw[otNavy,line width=1.2pt] plot[domain=-1:0.9,samples=100] ({\x},{\x*exp(\x)});
  \draw[otGreenDark,line width=1.2pt] plot[domain=-1:0.9,samples=100] ({\x*exp(\x)},{\x});
  \fill[otNavy] (-1,-0.3679) circle (1.4pt);
  \fill[otGreenDark] (-0.3679,-1) circle (1.4pt);
  \node[font=\scriptsize,text=otNavy,anchor=south west] at (0.75,1.9) {$y=x\,\mathrm{e}^{x}$};
  \node[font=\scriptsize,text=otGreenDark,anchor=west] at (2.05,0.75) {$W$};
  \node[font=\scriptsize,text=otMuted,anchor=north west] at (-0.36,-1) {$\bigl(-\tfrac1{\mathrm{e}},-1\bigr)$};
  \node[font=\tiny,text=otMuted,below] at (1,0) {$1$};
  \node[font=\tiny,text=otMuted,left] at (0,1) {$1$};
\end{tikzpicture}\\[0.5mm]
\legende{La fonction $x\mapsto x\,\mathrm{e}^{x}$ (en bleu) et sa réciproque
$W$ (en vert), symétriques par rapport à la première bissectrice ; $W$ a une
tangente verticale au point $\bigl(-\frac1{\mathrm{e}},-1\bigr)$.}
\end{center}

\partieExo{B}{Développement limité en $0$}
On admet que $W$ admet en $0$ un développement limité à tout ordre.
\begin{questions}[resume]
  \item En identifiant les coefficients dans $W(y)\,\mathrm{e}^{W(y)}=y$,
    montrer que $W(y)=y-y^{2}+\dfrac32y^{3}+o\bigl(y^{3}\bigr)$.
  \item Pousser le calcul à l'ordre $4$ ; on vérifiera que le coefficient
    de $y^{n}$ vaut $(-n)^{n-1}/n!$ pour $n\leqslant4$ (c'est vrai pour
    tout $n$ : formule d'inversion de Lagrange, admise).
  \item En déduire une valeur approchée de la solution de
    $x\,\mathrm{e}^{x}=0{,}01$, en précisant l'ordre de grandeur de
    l'erreur.
\end{questions}

\partieExo{C}{Comportement en $+\infty$}
\begin{questions}[resume]
  \item Montrer que $W(y)\to+\infty$ quand $y\to+\infty$, que
    $W(y)+\ln W(y)=\ln y$ pour $y>0$, et en déduire $W(y)\sim\ln y$.
  \item Établir successivement
    \[
    W(y)=\ln y-\ln\ln y+o(1),
    \qquad
    W(y)=\ln y-\ln\ln y+\frac{\ln\ln y}{\ln y}
    +o\Bigl(\frac{\ln\ln y}{\ln y}\Bigr).
    \]
  \item Pour $y>1$, montrer que l'équation $x^{x}=y$ admet une unique
    solution $x_y>1$, que
    \[
    x_y=\frac{\ln y}{W(\ln y)},
    \]
    et en déduire un équivalent simple de $x_y$ quand $y\to+\infty$.
\end{questions}
\end{defi}

